\documentclass[12pt, a4paper]{article}
\usepackage[T2A]{fontenc}
\usepackage[utf8]{inputenc}
\usepackage[russian]{babel}
\usepackage{amsmath, amssymb, amsthm}
\usepackage{geometry}
\geometry{left=2.5cm, right=2cm, top=2cm, bottom=2cm}

\theoremstyle{definition}
\newtheorem*{theory}{Теория}

\begin{document}

\begin{center}
    \Large\bfseries Задача №10: Решение нелинейных уравнений методом Чебышева \\
\end{center}

\vspace{0.5cm}

\section{Введение}
В данной задаче рассматривается численное решение нелинейных уравнений вида $f(x) = 0$ с использованием итерационного метода Чебышева высокого порядка. Этот метод обеспечивает более высокую скорость сходимости по сравнению со стандартным методом Ньютона за счет вычисления старших производных функции.

В работе реализован алгоритм для решения двух модельных функций:
\begin{enumerate}
    \item $f(x) = x^m - a = 0$ (вычисление корня $m$-й степени: $x^* = \sqrt[m]{a}$).
    \item $f(x) = e^x - a = 0$ (вычисление логарифма: $x^* = \ln(a)$.
\end{enumerate}

\section{Теоретическая база: Разложение обратной функции}
\begin{theory}
Пусть задано нелинейное уравнение $f(x) = 0$, и мы ищем его точный корень $z$, то есть $f(z) = 0$. 

Рассмотрим обратную функцию $x = F(y)$, по определению удовлетворяющую условию $F(f(x)) = x$. Если мы сможем найти значение обратной функции в нуле, мы точно найдем искомый корень: 
\[
z = F(0)
\]
Пусть $x_n$ --- текущее приближение к корню, а значение функции в этой точке равно $y_n = f(x_n)$. Разложим функцию $F(y)$ в ряд Тейлора в окрестности известной нам точки $y_n$, сохраняя члены вплоть до пятой производной (для обеспечения пятого порядка сходимости, используемого в программе):
\[
F(y) = F(y_n) + F'(y_n)(y - y_n) + \frac{F''(y_n)}{2!}(y - y_n)^2 + \dots + \frac{F^{(5)}(y_n)}{5!}(y - y_n)^5 + O((y-y_n)^6)
\]
Так как мы ищем $z = F(0)$, подставим $y = 0$ в это разложение. Учитывая, что $F(y_n) = x_n$, получаем итерационную формулу:
\[
x_{n+1} = x_n - F'(y_n)y_n + \frac{F''(y_n)}{2}y_n^2 - \frac{F'''(y_n)}{6}y_n^3 + \frac{F^{(4)}(y_n)}{24}y_n^4 - \frac{F^{(5)}(y_n)}{120}y_n^5
\]
Чтобы использовать эту формулу на практике, необходимо выразить производные обратной функции $F(y)$ через производные исходной функции $f(x)$, последовательно применяя правило дифференцирования сложной функции $\left(\frac{dx}{dy} = \frac{1}{f'(x)}\right)$:
\begin{align*}
F'(y) &= \frac{1}{f'(x)} \\
F''(y) &= \frac{d}{dy}\left(\frac{1}{f'(x)}\right) = \frac{dx}{dy} \cdot \frac{d}{dx}\left(\frac{1}{f'(x)}\right) = -\frac{f''(x)}{(f'(x))^3} \\
F'''(y) &= \frac{d}{dy}\left( -\frac{f''(x)}{(f'(x))^3} \right) = \frac{3(f''(x))^2 - f'(x)f'''(x)}{(f'(x))^5}
\end{align*}
\end{theory}

Уже на третьей производной формула становится громоздкой. Вычисление $F^{(4)}(y)$ и $F^{(5)}(y)$ в общем виде порождает длинные алгебраические дроби с комбинациями производных $f(x)$ вплоть до пятого порядка. 

Именно поэтому вычисление таких поправок «в лоб» на каждой итерации вычислительно крайне неэффективно. Однако для конкретных функций (таких как $x^m - a$ и $e^x - a$) можно обойти прямое вычисление этих производных, построив эквивалентное разложение аналитически.

\section{Вывод коэффициентов для модельных функций}
Для применения разложения на практике введем безразмерный параметр малости (нормированную невязку) $t$, зависящий от текущего шага $x_n$.

\subsection{Функция $f(x) = x^m - a = 0$}
Нам необходимо найти корень $z = \sqrt[m]{a} = a^{1/m}$. 
Введем параметр $t$, характеризующий относительное отклонение текущего приближения:
\[
t = 1 - \frac{a}{x_n^m} \implies \frac{a}{x_n^m} = 1 - t \implies a = x_n^m(1 - t)
\]
Подставим это выражение для $a$ в формулу точного корня:
\[
z = a^{1/m} = \left( x_n^m(1 - t) \right)^{1/m} = x_n (1 - t)^{1/m}
\]
Так как вблизи корня $x_n^m \approx a$, параметр $t \to 0$. Разложим выражение $(1 - t)^{1/m}$ в ряд Маклорена (обобщенный биномиальный ряд Ньютона) по степеням $t$:
\[
(1 - t)^\alpha = 1 - \alpha t + \frac{\alpha(\alpha - 1)}{2!} t^2 - \frac{\alpha(\alpha - 1)(\alpha - 2)}{3!} t^3 + \dots
\]
Полагая $\alpha = \frac{1}{m}$, последовательно вычислим коэффициенты $c_k$ при степенях $t^k$ вплоть до пятого порядка:
\begin{align*}
c_1 &= -\alpha = -\frac{1}{m} \\
c_2 &= \frac{\alpha(\alpha-1)}{2} = \frac{\frac{1}{m}\left(\frac{1}{m}-1\right)}{2} = \frac{1 - m}{2m^2} \\
c_3 &= -\frac{\alpha(\alpha-1)(\alpha-2)}{6} = -\frac{\frac{1}{m}\left(\frac{1-m}{m}\right)\left(\frac{1-2m}{m}\right)}{6} = -\frac{(1 - m)(1 - 2m)}{6m^3} \\
c_4 &= \frac{\alpha(\alpha-1)(\alpha-2)(\alpha-3)}{24} = \frac{(1 - m)(1 - 2m)(1 - 3m)}{24m^4} \\
c_5 &= -\frac{\alpha(\alpha-1)(\alpha-2)(\alpha-3)(\alpha-4)}{120} = -\frac{(1 - m)(1 - 2m)(1 - 3m)(1 - 4m)}{120m^5}
\end{align*}
Таким образом, итерационная формула принимает мультипликативный вид:
\[
x_{n+1} = x_n \left( 1 + c_1 t + c_2 t^2 + c_3 t^3 + c_4 t^4 + c_5 t^5 \right)
\]

\subsection{Функция $f(x) = e^x - a = 0$}
Нам необходимо найти корень $z = \ln(a)$. 
Аналогично введем параметр невязки $t$:
\[
t = 1 - \frac{a}{e^{x_n}} \implies \frac{a}{e^{x_n}} = 1 - t \implies a = e^{x_n}(1 - t)
\]
Подставим это выражение для $a$ в формулу точного корня, используя свойства логарифма:
\[
z = \ln(a) = \ln(e^{x_n}(1 - t)) = \ln(e^{x_n}) + \ln(1 - t) = x_n + \ln(1 - t)
\]
Разложим функцию $\ln(1 - t)$ в ряд Тейлора (ряд Меркатора), справедливый при $|t| < 1$:
\[
\ln(1 - t) = -t - \frac{t^2}{2} - \frac{t^3}{3} - \frac{t^4}{4} - \frac{t^5}{5} - \dots
\]
Коэффициенты разложения $c_k$ при степенях $t^k$ здесь являются строгими константами:
\[
c_1 = -1, \quad c_2 = -\frac{1}{2}, \quad c_3 = -\frac{1}{3}, \quad c_4 = -\frac{1}{4}, \quad c_5 = -\frac{1}{5}
\]
Итерационная формула принимает аддитивный вид:
\[
x_{n+1} = x_n + c_1 t + c_2 t^2 + c_3 t^3 + c_4 t^4 + c_5 t^5
\]

\section{Реализация алгоритма}
Алгоритм реализован на языке C++ с разбиением на интерфейс (\texttt{chebyshev.hpp}) и реализацию (\texttt{chebyshev.cpp}). Сборка осуществляется с помощью утилиты \texttt{make} со строгими флагами компиляции (\texttt{-O3 -Wall -Wextra -Werror}). Для оптимизации вычисления полинома 5-й степени в коде применена вычислительная \textbf{схема Горнера}, сокращающая количество ресурсоемких операций умножения.

\subsection{Критерии останова}
Итерационный процесс продолжается до тех пор, пока не выполнится одно из следующих условий:
\begin{enumerate}
    \item Достигнуто максимальное число итераций (\texttt{MAX\_ITER = 1000}).
    \item Малость приращения (сходимость по аргументу): $|x_{n+1} - x_n| < \varepsilon$.
    \item Малость невязки (сходимость по значению): $|f(x_{n+1})| < \varepsilon$.
\end{enumerate}
В программе задана точность $\varepsilon = 10^{-10}$, что обеспечивает нахождение корня с высокой достоверностью.

\subsection{Численный эксперимент}
Для проверки работоспособности алгоритма были проведены тесты на обеих модельных функциях. Программа продемонстрировала феноменальную скорость сходимости:
\begin{itemize}
    \item Для функции $f(x) = e^x - a$: алгоритм сошелся к корню $x^* \approx 1.3862943611$ всего за \textbf{7 итераций}. Итоговая невязка составила $\approx 8.88 \times 10^{-16}$.
    \item Для функции $f(x) = x^m - a$: алгоритм сошелся к корню $x^* \approx 1.2190136542$ за \textbf{8 итераций}. Итоговая невязка составила $\approx 1.33 \times 10^{-15}$.
\end{itemize}

\section{Выводы}
За счет предварительного аналитического вывода коэффициентов разложения (разложения обратной функции $F(0)=z$) мы можем блестяще адаптировать метод Чебышева для конкретных нелинейных функций, избегая дорогостоящего численного вычисления производных на каждом шаге итерации. 

Сохранение членов ряда вплоть до 5-го порядка ($t^5$) в программной реализации обеспечивает сверхбыструю сходимость. Как показал численный эксперимент, метод позволяет достигать \textbf{машинного нуля} (машинной точности $\varepsilon_{mach} \approx 10^{-16}$ для типа \texttt{double}) за единицы итераций. Использование метода Чебышева высокого порядка делает алгоритм сверхустойчивым и предпочтительным по сравнению со стандартным методом Ньютона при работе с экспоненциальными и степенными функциями.

\end{document}